\documentclass[10pt]{article}
\usepackage[margin=1in]{geometry}
\usepackage[T1]{fontenc}
\usepackage{lmodern,microtype,amsmath,amssymb,amsthm}
\theoremstyle{definition}
\newtheorem{assumption}{Assumption}
\begin{document}
\section*{Bounded Scalar Convex Optimization}
We consider minimizing a convex, $H$-Lipschitz function $f:\mathbb{R}\to\mathbb{R}$ on $C=[a,b]$, where $a\leq b$ and $H>0$. A given point $x^0\in C$ lies within distance $R\geq0$ of a minimizer $x^\star\in C$. The oracle returns a random scalar bounded in absolute value by $H$, whose conditional mean given past oracle randomness is a subgradient on $C$; replies may be dependent. For $\varepsilon>0$, the goal is a measurable, always feasible output $\bar{x}$ with integrable objective gap and $\mathbb{E}_{P}[f(\bar{x})-f(x^\star)]\leq\varepsilon$, where $P$ is the joint law of algorithm and oracle randomness.

The optimization problem is
\[
\min_{x\in C} f(x).
\]
The algorithm receives $a,b,H,R,\varepsilon$, and $x^0$, with the parameter ranges given above. Its randomness is independent of the oracle randomness.

Our assumptions:
\begin{assumption}[Convexity and Lipschitz continuity]
For every $x,y\in C$ and $t\in[0,1]$,
\[
\begin{aligned}
f((1-t)x+ty)&\leq(1-t)f(x)+tf(y),\\
|f(x)-f(y)|&\leq H|x-y|.
\end{aligned}
\]
\end{assumption}

\begin{assumption}[Initial distance]
There exists $x^\star\in C$ such that
\[
\begin{aligned}
f(x^\star)&\leq f(x) &&\text{for every }x\in C,\\
|x^0-x^\star|&\leq R.
\end{aligned}
\]
\end{assumption}

\begin{assumption}[Bounded conditional subgradients]
Let $\xi=(\xi_n)_{n\geq0}$ denote the oracle's random samples with joint law $\nu$, and let
\[
\mathcal{F}_n=\sigma(\xi_0,\ldots,\xi_{n-1}).
\]
At call $n$, the scalar reply has the form $G_n(h,x,\xi_n)$, where $h\in(\mathbb{R}\times\mathbb{R})^n$ is a history of query--reply pairs and $x\in\mathbb{R}$ is the query. For each $n$, the map $G_n$ is jointly measurable in its history, query, and sample arguments.

For every $\mathcal{F}_n$-measurable history $h_n$ and every $\mathcal{F}_n$-measurable query $X_n$ taking values in $C$ for every realization, define
\[
\begin{aligned}
g_n&=G_n(h_n,X_n,\xi_n),\\
s_n&=\mathbb{E}_{\nu}[g_n\mid\mathcal{F}_n].
\end{aligned}
\]
These replies satisfy
\[
\begin{aligned}
|g_n|&\leq H &&\nu\text{-a.s.},\\
f(X_n)+s_n(y-X_n)&\leq f(y)
&&\text{for all }y\in C,\quad\nu\text{-a.s.}
\end{aligned}
\]
The second inequality holds simultaneously for every $y\in C$ outside a single null set for each choice of $n,h_n,X_n$. The bound on $g_n$ ensures its integrability.
\end{assumption}

The oracle is defined for real queries; its boundedness and conditional subgradient guarantees apply to queries in $C$. After the allotted oracle calls, the output $\bar{x}$ must be measurable and belong to $C$ for every realization. Its objective gap must be integrable under $P$ and satisfy
\[
\mathbb{E}_{P}\bigl[f(\bar{x})-f(x^\star)\bigr]\leq\varepsilon.
\]
\end{document}
